\documentclass[aspectratio=169, 10pt]{beamer}

% --- PAQUETES BÁSICOS ---
\usepackage[utf8]{inputenc}
\usepackage[T1]{fontenc}
\usepackage[spanish, es-tabla]{babel}
\usepackage{amsmath}
\usepackage{amssymb}
\usepackage{booktabs}
\usepackage{tikz}
\usetikzlibrary{shapes.geometric, arrows, positioning}

% --- TEMA FIME UANL ---
\usetheme[showlogo, nosectionpage]{FIME}

% Estilos TikZ adaptados
\tikzset{
    block/.style={rectangle, draw=FIMEgreen, fill=FIMElightgray, text width=2.4cm, text centered, rounded corners, minimum height=1.2cm, font=\scriptsize\bfseries, line width=1pt},
    arrow/.style={thick,->,>=stealth,draw=UANLblue,line width=1.2pt},
    actor/.style={rectangle, draw=UANLblue, fill=FIMElightgray, text width=2.8cm, text centered, rounded corners, minimum height=1cm, font=\small\bfseries, line width=1pt}
}

% --- METADATOS ---
\title[Herramientas Multicriterio - Sesión 2]{Herramientas Multicriterio para Toma de Decisiones}
\subtitle{Sesión 2: Ponderación Subjetiva: AHP (Saaty) y Best-Worst Method (BWM)}
\author[Leonardo G. Hernández]{Dr. Leonardo Gabriel Hernández Landa}
\institute[FIME - UANL]{
    Doctorado en Logística y Cadena de Suministro\\
    Facultad de Ingeniería Mecánica y Eléctrica (FIME)\\
    Universidad Autónoma de Nuevo León (UANL)
}
\date{Tetramestre de Otoño}

% --- INICIO DEL DOCUMENTO ---
\begin{document}

% Slide 1: Portada
\begin{frame}[plain]
    \titlepage
\end{frame}

% Slide 2: Objetivos de la Sesión
\begin{frame}{Objetivos de la Sesión}
    \begin{block}{Al finalizar esta sesión, el estudiante de doctorado será capaz de:}
        \begin{itemize}
            \item \textbf{Analizar críticamente} los axiomas matemáticos del Proceso de Jerarquía Analítica (AHP) de Thomas Saaty.
            \item \textbf{Evaluar} la validez científica y el fenómeno de \textit{Rank Reversal} en AHP (Debate Dyer-Saaty).
            \item \textbf{Dominar el algoritmo} y la formulación de optimización minimax del Best-Worst Method (BWM) desarrollado por Jafar Rezaei.
            \item \textbf{Modelar y resolver} problemas de ponderación subjetiva aplicados a la toma de decisiones estratégicas en logística y cadena de suministro.
        \end{itemize}
    \end{block}
\end{frame}

% Slide 3: Agenda de la Sesión
\begin{frame}{Agenda de la Sesión (3 Horas)}
    \begin{columns}[T]
        \begin{column}{0.31\textwidth}
            \begin{block}{Bloque 1 (90 min)}
                \textbf{Fundamentos y Axiomas}
                \begin{itemize}
                    \scriptsize
                    \item Axiomas del AHP
                    \item Escala de Saaty (1-9)
                    \item Vector de prioridad y consistencia
                    \item El Debate sobre Rank Reversal
                    \item Introducción a BWM (Rezaei 2015)
                \end{itemize}
            \end{block}
        \end{column}
        
        \begin{column}{0.31\textwidth}
            \begin{exampleblock}{Bloque 2 (45 min)}
                \textbf{Debate de Literatura}
                \begin{itemize}
                    \scriptsize
                    \item Dyer (1990) vs. Saaty (1990)
                    \item Rezaei (2015)
                    \item Axiomática de comparación pareada
                    \item Carga cognitiva del decisor en logística
                \end{itemize}
            \end{exampleblock}
        \end{column}
        
        \begin{column}{0.31\textwidth}
            \begin{alertblock}{Bloque 3 (45 min)}
                \textbf{Taller Práctico}
                \begin{itemize}
                    \scriptsize
                    \item Cálculo de consistencia en Excel
                    \item Formulación lineal de BWM
                    \item Resolución con Gurobi/SciPy
                    \item Selección de transportista LTL
                \end{itemize}
            \end{alertblock}
        \end{column}
    \end{columns}
\end{frame}

% Slide 4: Introducción a la Ponderación Subjetiva
\begin{frame}{Ponderación de Criterios en Decisiones Multicriterio}
    \begin{itemize}
        \item La importancia relativa de los criterios ($w_j$) es un elemento central de los modelos MADM/MCDM.
        \item \textbf{Ponderación Subjetiva:} Se basa enteramente en las preferencias de expertos u tomadores de decisión (DM).
        \item \textbf{Ponderación Objetiva:} Se deduce analíticamente de la dispersión de datos en la matriz de decisión (ej. Entropía, CRITIC).
    \end{itemize}
    \begin{block}{Desafío en el Nivel Doctoral}
        El estudiante debe justificar la representatividad matemática de las prioridades elicitaedas:
        \begin{itemize}
            \item ¿Cómo capturar la racionalidad humana sin inducir sesgos cognitivos o fatiga en el decisor?
            \item ¿Qué supuestos axiomáticos rigen a cada método?
        \end{itemize}
    \end{block}
\end{frame}

% Slide 5: Fundamentos del Proceso de Jerarquía Analítica (AHP)
\begin{frame}{El Proceso de Jerarquía Analítica (AHP)}
    \begin{columns}[T]
        \begin{column}{0.6\textwidth}
            \begin{itemize}
                \item Desarrollado por Thomas L. Saaty en la década de 1970.
                \item Descompone un problema complejo en una estructura jerárquica:
                      \textbf{Meta $\rightarrow$ Criterios $\rightarrow$ Subcriterios $\rightarrow$ Alternativas}.
                \item Reemplaza la asignación directa de pesos por \textbf{comparaciones pareadas} sistemáticas.
                \item Genera una matriz de comparación recíproca positiva $\mathbf{A} = [a_{ij}]$.
            \end{itemize}
        \end{column}
        \begin{column}{0.38\textwidth}
            \begin{exampleblock}{Ventajas Metodológicas}
                \scriptsize
                \begin{itemize}
                    \item Permite estructurar juicios tanto cualitativos como cuantitativos.
                    \item Proporciona una medida formal de consistencia para los juicios del decisor.
                    \item Es el método MCDM más difundido y validado en la industria global.
                \end{itemize}
            \end{exampleblock}
        \end{column}
    \end{columns}
\end{frame}

% Slide 6: Los Cuatro Axiomas de Saaty
\begin{frame}{Cimiento Axiomático del AHP}
    La validez matemática del AHP descansa sobre cuatro axiomas\footnote{\textbf{Axioma:} Principio o regla básica que se asume como verdadera sin demostración para servir de punto de partida.} estrictos (Saaty, 1986):
    \begin{enumerate}
        \item \textbf{Reciprocidad:} Si el criterio $i$ es $x$ veces más importante que el criterio $j$, entonces el criterio $j$ es $1/x$ veces más importante que el criterio $i$ ($a_{ji} = 1/a_{ij}$).
        \item \textbf{Homogeneidad:} Los elementos que se comparan deben pertenecer a la misma escala de magnitud (es decir, no se debe comparar elementos infinitamente mayores con otros infinitamente pequeños).
        \item \textbf{Dependencia:} Las prioridades de los elementos de un nivel de la jerarquía dependen de las del nivel superior inmediato.
        \item \textbf{Expectativas:} Cualquier cambio en la estructura de la jerarquía requiere revaluar las expectativas del tomador de decisiones.
    \end{enumerate}
\end{frame}

% Slide 7: La Escala de Saaty (1 a 9)
\begin{frame}{La Escala Fundamental de Comparación}
    Saaty propuso una escala de medición discreta basada en estudios de psicología cognitiva (Ley de Miller de la capacidad humana de procesamiento de información, $7 \pm 2$):
    \begin{table}
        \scriptsize
        \begin{tabular}{cll}
            \toprule
            \textbf{Intensidad} & \textbf{Definición Verbal} & \textbf{Explicación Lógica} \\
            \midrule
            1 & Igual importancia & Dos actividades contribuyen igual a la meta. \\
            3 & Importancia moderada & La experiencia/juicio favorece ligeramente a una sobre otra. \\
            5 & Importancia fuerte & La experiencia/juicio favorece fuertemente a una sobre otra. \\
            7 & Importancia muy fuerte o demostrada & Una actividad es favorecida muy fuertemente y es dominante. \\
            9 & Importancia extrema & La evidencia que favorece a una actividad es del orden más alto. \\
            2, 4, 6, 8 & Valores intermedios & Utilizados cuando se busca un compromiso entre juicios. \\
            \bottomrule
        \end{tabular}
    \end{table}
    \begin{alertblock}{Crítica Axiomática}
        ¿Es la escala lineal de 1-9 la mejor representación del juicio humano? Escalas geométricas, exponenciales o logarítmicas son frecuentemente utilizadas en investigación avanzada.
    \end{alertblock}
\end{frame}

% Slide 8: Construcción de la Matriz AHP
\begin{frame}{Construcción de la Matriz de Comparaciones Pareadas}
    Sea un conjunto de criterios $\{C_1, C_2, \dots, C_n\}$, la matriz de comparación $\mathbf{A} \in \mathbb{R}^{n \times n}$ se define como:
    $$\mathbf{A} = \begin{pmatrix}
        1 & a_{12} & \dots & a_{1n} \\
        1/a_{12} & 1 & \dots & a_{2n} \\
        \vdots & \vdots & \ddots & \vdots \\
        1/a_{1n} & 1/a_{2n} & \dots & 1
    \end{pmatrix}$$
    \begin{itemize}
        \item Donde $a_{ij} > 0$ representa la preferencia del criterio $C_i$ sobre $C_j$.
        \item El número total de comparaciones requeridas para $n$ criterios es:
              $$N = \frac{n(n-1)}{2}$$
        \item Para $n=10$ criterios, ¡se requieren 45 comparaciones! Riesgo de fatiga e inconsistencia en comités logísticos.
    \end{itemize}
\end{frame}

% Slide 9: Vector de Prioridades y el Autovector Principal
\begin{frame}{Elicitación\footnote{\textbf{Elicitación (de prioridades):} Proceso sistemático para extraer y formalizar conocimientos o juicios subjetivos de un experto.} del Vector de Prioridades}
    Si el decisor fuera perfectamente consistente, se cumpliría que $a_{ij} = w_i / w_j$, lo que implica que:
    $$a_{ij} \frac{w_j}{w_i} = 1 \implies \sum_{j=1}^n a_{ij} w_j = n w_i \implies \mathbf{A}\mathbf{w} = n\mathbf{w}$$
    Dado que el ser humano comete imprecisiones, los juicios tienen ruido. Saaty demostró que para una matriz recíproca positiva general:
    $$\mathbf{A}\mathbf{w} = \lambda_{\max}\mathbf{w}$$
    \begin{block}{Método del Autovector\footnote{\textbf{Autovector (Eigenvector):} Vector especial que no cambia de dirección bajo una transformación. En AHP, representa el peso relativo estable de los criterios.} Principal (Principal Eigenvector Method)}
        \begin{itemize}
            \item Se calcula el vector propio $\mathbf{w}$ correspondiente al autovalor máximo $\lambda_{\max}$ de la matriz $\mathbf{A}$.
            \item Teorema de Perron-Frobenius garantiza que existe un único autovector positivo para una matriz positiva irreductible.
            \item Siempre se cumple que $\lambda_{\max} \ge n$, alcanzando la igualdad sólo si la matriz es perfectamente consistente.
        \end{itemize}
    \end{block}
\end{frame}

% Slide 10: Medición de la Consistencia en AHP
\begin{frame}{El Índice y la Razón de Consistencia}
    Para cuantificar la desviación de la consistencia perfecta, Saaty formuló:
    \begin{columns}[T]
        \begin{column}{0.48\textwidth}
            \begin{block}{Índice de Consistencia (CI)}
                $$CI = \frac{\lambda_{\max} - n}{n-1}$$
            \end{block}
        \end{column}
        \begin{column}{0.48\textwidth}
            \begin{block}{Razón de Consistencia (CR)}
                $$CR = \frac{CI}{RI}$$
            \end{block}
        \end{column}
    \end{columns}
    \vspace{0.2cm}
    \begin{itemize}
        \item Donde $RI$ es el \textbf{Índice Aleatorio} (consistencia promedio de matrices generadas al azar de tamaño $n$):
    \end{itemize}
    \begin{table}
        \scriptsize
        \begin{tabular}{cccccccccc}
            \toprule
            $n$ & 1 & 2 & 3 & 4 & 5 & 6 & 7 & 8 & 9 \\
            \midrule
            $RI$ & 0.00 & 0.00 & 0.58 & 0.90 & 1.12 & 1.24 & 1.32 & 1.41 & 1.45 \\
            \bottomrule
        \end{tabular}
    \end{table}
    \begin{alertblock}{Regla de Decisión}
        Si $CR < 0.10$ (10\%), los juicios son lo suficientemente consistentes para ser aceptados. De lo contrario, el decisor debe revisar y reformular la matriz.
    \end{alertblock}
\end{frame}

% Slide 11: El Gran Debate: Rank Reversal (Inversión de Rangos)
\begin{frame}{La Polémica de Dyer vs. Saaty}
    \begin{columns}[T]
        \begin{column}{0.5\textwidth}
            \begin{block}{El Fenómeno de Rank Reversal}
                La adición o eliminación de una alternativa \textit{irrelevante} (o duplicada) en la matriz de decisión puede provocar que el ranking relativo de las alternativas originales se invierta.
            \end{block}
            \begin{itemize}
                \item \textbf{Dyer (1990):} Argumentó que esto viola el axioma de independencia de alternativas irrelevantes de Luce y Arrow, invalidando al AHP clásico.
            \end{itemize}
        \end{column}
        \begin{column}{0.48\textwidth}
            \begin{block}{La Respuesta de Saaty (1990)}
                Defendió que la inversión de rangos es un comportamiento natural de la mente humana que ajusta preferencias relativas al introducir nuevos elementos en el contexto (ej. el efecto señuelo en marketing).
            \end{block}
            \begin{exampleblock}{Modos de AHP}
                \scriptsize
                \begin{itemize}
                    \item \textbf{Distributivo:} Normalización clásica (permite Rank Reversal).
                    \item \textbf{Ideal:} Normaliza dividiendo por la mejor alternativa (evita Rank Reversal).
                \end{itemize}
            \end{exampleblock}
        \end{column}
    \end{columns}
\end{frame}

% Slide 12: Limitaciones de AHP
\begin{frame}{Debilidades Operativas del AHP Clásico}
    \begin{itemize}
        \item \textbf{Alta Carga Cognitiva:} Requiere $n(n-1)/2$ comparaciones pareadas. Para problemas grandes con múltiples criterios y alternativas, el número de preguntas se vuelve inviable.
        \item \textbf{Inconsistencia Estructural:} A mayor número de comparaciones, mayor es la inconsistencia natural generada por el decisor humano, lo que obliga a repetir encuestas.
        \item \textbf{Escala Asimétrica:} La escala de Saaty 1-9 no es lineal respecto a los pesos finales, lo que genera distorsiones si no se normaliza adecuadamente.
    \end{itemize}
    \begin{block}{Pregunta de Investigación}
        ¿Es posible obtener pesos subjetivos robustos con menor cantidad de comparaciones pareadas y garantizando consistencia matemática?
    \end{block}
\end{frame}

% Slide 13: Introducción al Best-Worst Method (BWM)
\begin{frame}{Best-Worst Method (BWM) de Rezaei (2015)}
    \begin{columns}[T]
        \begin{column}{0.58\textwidth}
            \begin{itemize}
                \item Propuesto por Jafar Rezaei en 2015 (\textit{Omega}).
                \item Es un método de comparación pareada estructurado y simplificado.
                \item En lugar de construir una matriz completa, el decisor solo compara:
                      \begin{enumerate}
                          \item El mejor criterio respecto a todos los demás (\textbf{Best-to-Others}).
                          \item Todos los criterios respecto al peor (\textbf{Others-to-Worst}).
                      \end{enumerate}
                \item Requiere únicamente:
                      $$N_{BWM} = 2n - 3$$
                      comparaciones pareadas.
            \end{itemize}
        \end{column}
        \begin{column}{0.38\textwidth}
            \begin{exampleblock}{Ventajas de BWM}
                \scriptsize
                \begin{itemize}
                    \item Menor carga cognitiva.
                    \item Juicios significativamente más consistentes.
                    \item Formulación matemática lineal y convexa\footnote{\textbf{Convexidad:} Propiedad matemática que asegura que cualquier óptimo local es el óptimo global (el mejor absoluto), facilitando su resolución.} que garantiza una solución única global.
                \end{itemize}
            \end{exampleblock}
        \end{column}
    \end{columns}
\end{frame}

% Slide 14: Los Pasos del Algoritmo BWM
\begin{frame}{Algoritmo de BWM: Comparaciones Estructuradas}
    \begin{enumerate}
        \item \textbf{Paso 1:} Definir el conjunto de criterios logísticos $\{C_1, C_2, \dots, C_n\}$.
        \item \textbf{Paso 2:} Identificar el mejor criterio ($C_B$) y el peor criterio ($C_W$).
        \item \textbf{Paso 3:} Elicitar el vector de comparación del mejor criterio respecto a los demás (\textbf{Best-to-Others}, escala 1-9):
              $$A_B = (a_{B1}, a_{B2}, \dots, a_{Bn})$$
              *(Donde $a_{BB} = 1$)*
        \item \textbf{Paso 4:} Elicitar el vector de comparación de todos los criterios respecto al peor (\textbf{Others-to-Worst}, escala 1-9):
              $$A_W = (a_{1W}, a_{2W}, \dots, a_{nW})^T$$
              *(Donde $a_{WW} = 1$)*
    \end{enumerate}
\end{frame}

% Slide 15: Formulación Matemática Minimax de BWM
\begin{frame}{Formulación del Modelo Minimax\footnote{\textbf{Minimax:} Enfoque de optimización que busca minimizar la máxima desviación o peor error en un sistema de ecuaciones.} No Lineal}
    El objetivo es determinar los pesos óptimos $(w_1^*, w_2^*, \dots, w_n^*)$ que minimicen la desviación absoluta máxima de las razones de pesos respecto a las comparaciones del decisor:
    \begin{columns}[T]
        \begin{column}{0.48\textwidth}
            $$\min \xi$$
            $$\text{s.t.}$$
            $$\left| \frac{w_B}{w_j} - a_{Bj} \right| \le \xi, \quad \forall j$$
            $$\left| \frac{w_j}{w_W} - a_{jW} \right| \le \xi, \quad \forall j$$
            $$\sum_{j=1}^n w_j = 1, \quad w_j \ge 0, \quad \forall j$$
        \end{column}
        \begin{column}{0.48\textwidth}
            \begin{alertblock}{Propiedades del Modelo}
                \scriptsize
                \begin{itemize}
                    \item Este problema no es lineal y puede poseer múltiples soluciones locales u óptimas no convexas.
                    \item Para mitigar esto, Rezaei (2016) reformuló el problema de manera lineal para evitar denominadores variables.
                \end{itemize}
            \end{alertblock}
        \end{column}
    \end{columns}
\end{frame}

% Slide 16: Modelo BWM Lineal de Rezaei (2016)
\begin{frame}{Formulación Matemática del BWM Linealizado}
    Multiplicando las restricciones por sus denominadores, Rezaei transformó el modelo minimax en un \textbf{Problema de Programación Lineal (PPL)} convexo con solución única y global:
    \begin{block}{Modelo BWM Lineal}
        $$\min \xi^L$$
        $$\text{s.t.}$$
        $$\left| w_B - a_{Bj} w_j \right| \le \xi^L, \quad \forall j$$
        $$\left| w_j - a_{jW} w_W \right| \le \xi^L, \quad \forall j$$
        $$\sum_{j=1}^n w_j = 1, \quad w_j \ge 0, \quad \forall j$$
    \end{block}
    \begin{exampleblock}{Resolución en Python}
        \scriptsize
        Se puede programar en un bloque simple utilizando la librería `scipy.optimize.linprog` o a través de `PuLP`.
    \end{exampleblock}
\end{frame}

% Slide 17: Consistencia en el Método BWM
\begin{frame}{Medición del Grado de Consistencia en BWM}
    El valor óptimo resultante del modelo $\xi^*$ (o $\xi^L$ en el modelo lineal) indica el nivel de inconsistencia de los juicios pareados. La \textbf{Razón de Consistencia (CR)} se calcula como:
    $$CR = \frac{\xi^*}{\text{Consistencia Mínima (CI)}}$$
    \begin{table}
        \scriptsize
        \begin{tabular}{cccccccccc}
            \toprule
            $a_{BW}$ (Escala) & 1 & 2 & 3 & 4 & 5 & 6 & 7 & 8 & 9 \\
            \midrule
            $CI$ (Max No-lineal) & 0.00 & 0.44 & 1.00 & 1.63 & 2.30 & 3.00 & 3.73 & 4.47 & 5.23 \\
            \bottomrule
        \end{tabular}
    \end{table}
    \begin{itemize}
        \item A mayor valor de $a_{BW}$ (diferencia extrema entre el mejor y peor criterio), mayor es el error admisible $\xi^*$ en la escala de Saaty.
        \item Valores de $CR$ cercanos a 0 indican alta consistencia; valores superiores a 0.10 requieren revisión de los vectores de entrada.
    \end{itemize}
\end{frame}

% Slide 18: Comparación Directa: AHP vs. BWM
\begin{frame}{Comparativa Metodológica para Proyectos Doctorales}
    \begin{table}
        \scriptsize
        \begin{tabular}{lll}
            \toprule
            \textbf{Dimensión} & \textbf{Analytic Hierarchy Process (AHP)} & \textbf{Best-Worst Method (BWM)} \\
            \midrule
            \textbf{Complejidad} & Alta: $n(n-1)/2$ comparaciones & Baja: $2n-3$ comparaciones \\
            \textbf{Consistencia} & Propensa a errores de transitividad\footnote{\textbf{Transitividad:} Principio de consistencia lógica: si se prefiere A a B y B a C, lógicamente se debe preferir A a C. Los decisores humanos suelen violarlo por fatiga o confusión.} & Robusta debido a consistencia lógica \\
            \textbf{Optimización} & No requiere (Autovectores) & Requiere resolver un PPL minimax \\
            \textbf{Juicios de Entrada} & Matriz completa $n \times n$ & Dos vectores estructurados (1 $\times n$ y $n \times 1$) \\
            \textbf{Carga Cognitiva} & Exponencial con el tamaño del problema & Lineal con el tamaño del problema \\
            \textbf{Escala de Comparación} & Recíproca lineal (1-9) & Recíproca lineal (1-9) \\
            \bottomrule
        \end{tabular}
    \end{table}
    \begin{block}{Criterio de Selección para el Investigador}
        Para problemas logísticos con más de 5 criterios o con comités directivos de agenda limitada, BWM ofrece resultados equivalentes con un 60\% menos de esfuerzo computacional y cognitivo.
    \end{block}
\end{frame}

% Slide 19: Caso de Estudio: Selección de Transportista LTL
\begin{frame}{Selección de Transportista Terrestre Consolidado (LTL)}
    \begin{columns}[T]
        \begin{column}{0.48\textwidth}
            Un CEDIS de distribución de e-commerce en Monterrey debe ponderar 4 criterios operativos clave para subcontratar transportistas LTL:
            \begin{itemize}
                \item \textbf{Costo ($C_1$):} Costo de flete por kilómetro.
                \item \textbf{Confiabilidad ($C_2$):} Exactitud de entrega (On-Time In-Full).
                \item \textbf{Visibilidad ($C_3$):} Capacidad de rastreo GPS en tiempo real.
                \item \textbf{Sustentabilidad ($C_4$):} Emisiones promedio y flota ecológica.
            \end{itemize}
        \end{column}
        
        \begin{column}{0.48\textwidth}
            \begin{exampleblock}{Parámetros de Entrada de BWM}
                \scriptsize
                \begin{itemize}
                    \item \textbf{Criterio Best elegido:} Confiabilidad ($C_2$).
                    \item \textbf{Criterio Worst elegido:} Sustentabilidad ($C_4$).
                    \item \textbf{Vector Best-to-Others ($A_B$):}
                          $$a_{21}=2, \quad a_{22}=1, \quad a_{23}=3, \quad a_{24}=6$$
                    \item \textbf{Vector Others-to-Worst ($A_W$):}
                          $$a_{14}=3, \dots, a_{24}=6, \dots, a_{34}=2, \dots, a_{44}=1$$
                \end{itemize}
            \end{exampleblock}
        \end{column}
    \end{columns}
\end{frame}

% Slide 20: Taller Práctico y Trabajo Autónomo
\begin{frame}{Taller Aplicado del Caso Práctico 1 (E1)}
    \begin{itemize}
        \item El objetivo del **Caso Práctico 1 (E1)** es comparar empírica y estadísticamente los pesos resultantes de AHP y BWM para la selección del transportista LTL.
    \end{itemize}
    \begin{block}{Pauta de Trabajo para el Estudiante}
        \begin{enumerate}
            \item Construir la matriz AHP de comparaciones pareadas consistente con las preferencias de BWM.
            \item Calcular el vector propio y validar la consistencia ($CR < 10\%$).
            \item Implementar la formulación matemática del BWM lineal en Python (utilizando `PuLP` o `scipy.optimize.linprog`).
            \item Resolver el PPL, reportar los pesos $w^*$ y la inconsistencia óptima $\xi^L$.
            \item Contrastar la estabilidad de los rankings utilizando coeficientes de correlación y discutir el impacto de la fatiga del decisor en cadena de suministro.
        \end{enumerate}
    \end{block}
\end{frame}

% Slide 21: Conclusiones
\begin{frame}{Conclusiones Clave de la Ponderación Subjetiva}
    \begin{itemize}
        \item La elección de un método de ponderación subjetiva no es trivial; impacta directamente en la estabilidad del ranking en la cadena de suministro.
        \item AHP es el estándar comercial de la industria, pero sus supuestos axiomáticos de consistencia obligan a simplificar las jerarquías.
        \item BWM es una extensión metodológica del estado del arte que reduce la complejidad matemática y optimiza la precisión en la toma de decisiones grupales.
        \item Ambos métodos son pilares fundamentales para el **Producto Integrador de Aprendizaje (PIA)** de este curso doctoral.
    \end{itemize}
\end{frame}

% Slide 22: Referencias Bibliográficas Clave
\begin{frame}{Referencias Obligatorias y de Debate}
    \begin{itemize}
        \scriptsize
        \item \textbf{Saaty, T. L. (1980).} \textit{The Analytic Hierarchy Process}. McGraw-Hill.
        \item \textbf{Rezaei, J. (2015).} Best-worst multi-criteria decision-making method. \textit{Omega}, 53, 49--57.
        \item \textbf{Rezaei, J. (2016).} Non-linear and linear models in best-worst multiple criteria decision-making. \textit{Computers \& Industrial Engineering}, 96, 19--25.
        \item \textbf{Dyer, J. S. (1990).} Remarks on the analytic hierarchy process. \textit{Management Science}, 36(3), 249--258.
        \item \textbf{Saaty, T. L. (1990).} An exposition of the AHP in reply to the paper "Remarks on the analytic hierarchy process". \textit{Management Science}, 36(3), 259--268.
    \end{itemize}
\end{frame}

\end{document}
